\documentclass[11pt]{article}

\usepackage[margin=1in]{geometry}
\usepackage[T1]{fontenc}
\usepackage[utf8]{inputenc}
\usepackage{lmodern}
\usepackage{amsmath, amssymb, amsfonts}
\usepackage{mathtools}
\usepackage{enumitem}
\usepackage{hyperref}

\title{Mathematical Definitions of the Color-Cue Stimuli}
\author{}
\date{}

\begin{document}

\maketitle

This document is a compact mathematical companion to the \texttt{color\_cue}
workflow. Its goal is to describe, in equations, the \textbf{old stimulus}
used in the original color-cue exploration and the \textbf{new stimulus} used
in the theta-noise psychophysics design.

We keep the notation close to the implementation in the package modules
\texttt{color\_cue.stimulus} and \texttt{color\_cue.psychophysics}.

\section{Coordinate System and Cue-Angle Mapping}

The stimulus is defined on a two-dimensional image lattice with pixel
coordinates
\[
(u, v) \in \{1, \dots, H\} \times \{1, \dots, W\}.
\]

The latent color cue is parameterized by a scalar \texttt{cue} in $[0,1]$,
which is mapped to a hue angle $\theta$ on a restricted arc of the color
circle:
\[
\theta = \theta_{\min} + \text{cue}\,(\theta_{\max} - \theta_{\min}).
\]

Equivalently,
\[
\text{cue} = \frac{\theta - \theta_{\min}}{\theta_{\max} - \theta_{\min}}.
\]

In the current implementation,
\[
\theta_{\min} = -\frac{\pi}{2},
\qquad
\theta_{\max} = 0.
\]

So the cue lives on the quarter-circle from blue-ish angles toward red-ish
angles.

\section{Old Stimulus: Original Color-Cue Generator}

The original generator creates a spatially correlated complex random field and
then adds a deterministic signal vector of magnitude $\kappa$ pointing at angle
$\theta$.

\subsection{Frequency-Domain Noise}

Let $(f_x, f_y)$ denote spatial frequencies. The code defines the radial
frequency
\[
f = \sqrt{f_x^2 + f_y^2},
\]

and a power-law spatial spectrum
\[
S(f) = (f + \varepsilon)^{\beta},
\]

with a small $\varepsilon > 0$ for numerical stability and with the DC term
set to zero:
\[
S(0) = 0.
\]

A complex Gaussian field is sampled in the Fourier domain:
\[
\tilde{N}(f_x, f_y)
= \bigl(Z_1(f_x, f_y) + i Z_2(f_x, f_y)\bigr)\,S(f)^{1/2},
\]

where $Z_1, Z_2 \sim \mathcal{N}(0,1)$ independently.

The corresponding spatial complex noise field is then
\[
N(u,v) = \mathcal{F}^{-1}\!\left[\tilde{N}(f_x, f_y)\right](u,v).
\]

\subsection{Deterministic Color Signal}

The latent cue angle $\theta$ is converted into a unit complex vector
\[
e^{i\theta} = \cos\theta + i\sin\theta.
\]

The stimulus reliability parameter $\kappa \ge 0$ sets the magnitude of the
deterministic signal contribution. The complex image before taking angles is
\[
X(u,v) = N(u,v) + \kappa e^{i\theta}.
\]

Interpretation:
\begin{itemize}[leftmargin=2em]
\item $N(u,v)$ provides spatially correlated texture-like color fluctuations.
\item $\kappa e^{i\theta}$ shifts every pixel in the same direction in the
complex plane.
\item Larger $\kappa$ means the mean color signal dominates the texture noise
more strongly.
\end{itemize}

\subsection{Pixel-Wise Angle Readout}

Each pixel is then mapped to its complex angle:
\[
\phi(u,v) = \arg\bigl(X(u,v)\bigr).
\]

Finally, this angle is rescaled into a value in $[0,1)$ for lookup in the
color map:
\[
C(u,v) = \frac{1}{2\pi}\phi(u,v) + \mathbf{1}\{\phi(u,v) \le 0\}.
\]

The resulting array $C(u,v)$ is what is displayed using the custom colormap.

This is the \textbf{old single-stimulus definition}: one latent angle
$\theta$, one reliability $\kappa$, one correlated texture realization.

\section{Old Stimulus With Temporal Correlation}

For movies, the generator introduces temporal correlation in the Fourier-domain
noise.

If $\tilde{N}_t$ is the complex Fourier noise at frame $t$, then for
$t \ge 2$ the update is
\[
\tilde{N}_t
= \Gamma \odot \tilde{N}_{t-1} + (1-\Gamma) \odot \tilde{\eta}_t,
\]

where:
\begin{itemize}[leftmargin=2em]
\item $\tilde{\eta}_t$ is a fresh complex Gaussian draw with the same spectrum,
\item $\Gamma = e^{-\alpha f}$ is a frequency-dependent temporal persistence
term,
\item $\odot$ denotes elementwise multiplication.
\end{itemize}

Each frame then uses the same signal term $\kappa e^{i\theta}$ but a
temporally correlated noise realization.

\section{Why the Old Stimulus Motivates a New Noise Manipulation}

In the original exploration, two summary statistics behaved very differently:

\begin{enumerate}[leftmargin=2em]
\item \textbf{Arithmetic mean of pixel angles}
\[
\bar{\phi}_{\mathrm{arith}} = \frac{1}{HW} \sum_{u,v} \phi(u,v),
\]
which showed noticeable variability.

\item \textbf{Circular mean angle}
\[
\bar{\phi}_{\mathrm{circ}}
= \arg\left(\frac{1}{HW} \sum_{u,v} e^{i\phi(u,v)}\right),
\]
which showed little or no variability under the tested conditions.
\end{enumerate}

This suggests that if an observer behaves more like a pooling mechanism near
the circular mean, then much of the old texture noise may average away. That
is the motivation for injecting noise directly into the latent signal angle
itself.

\section{New Stimulus: Theta-Noise Comparison Design}

The new design is not a single stimulus in isolation. It is a
\textbf{pairwise comparison stimulus} used in a two-alternative forced-choice
(2AFC) task.

Each trial starts from:
\begin{itemize}[leftmargin=2em]
\item a base angle $\theta_0$,
\item a signed intended separation $\Delta\theta$,
\item an external noise scale $\sigma_{\mathrm{ext}}$.
\end{itemize}

The noiseless target angles for the left and right stimuli are defined
symmetrically:
\[
\theta_L^{\ast} = \theta_0 - \frac{\Delta\theta}{2},
\qquad
\theta_R^{\ast} = \theta_0 + \frac{\Delta\theta}{2}.
\]

So the intended left-right difference is exactly
\[
\theta_R^{\ast} - \theta_L^{\ast} = \Delta\theta.
\]

\subsection{Independent External Theta Noise}

In the main psychophysics design, each side receives its own independent
perturbation:
\[
\varepsilon_L \sim \mathcal{N}(0, \sigma_{\mathrm{ext}}^2),
\qquad
\varepsilon_R \sim \mathcal{N}(0, \sigma_{\mathrm{ext}}^2),
\]

with $\varepsilon_L$ and $\varepsilon_R$ independent.

The displayed latent angles become
\[
\theta_L = \theta_L^{\ast} + \varepsilon_L,
\qquad
\theta_R = \theta_R^{\ast} + \varepsilon_R.
\]

After this, each side is rendered using the \textbf{same old single-stimulus
generator}, but with its own noisy angle:
\[
X_L(u,v) = N_L(u,v) + \kappa e^{i\theta_L},
\qquad
X_R(u,v) = N_R(u,v) + \kappa e^{i\theta_R}.
\]

So the new design keeps the pixel generator structurally the same while
changing the trial-level latent parameterization.

\subsection{Shared External Theta Noise Control}

The control condition uses a single common perturbation
\[
\varepsilon \sim \mathcal{N}(0, \sigma_{\mathrm{ext}}^2),
\]

and applies it to both sides:
\[
\theta_L = \theta_L^{\ast} + \varepsilon,
\qquad
\theta_R = \theta_R^{\ast} + \varepsilon.
\]

This changes the absolute hue of the whole trial, but preserves the intended
relative difference:
\[
\theta_R - \theta_L
= (\theta_R^{\ast} + \varepsilon) - (\theta_L^{\ast} + \varepsilon)
= \Delta\theta.
\]

That is why shared noise is expected to affect discrimination much less than
independent noise.

\section{Task Definition in the New Experiment}

On each trial, the observer sees the left and right stimuli side by side and
answers one of two prompts:
\begin{itemize}[leftmargin=2em]
\item \textbf{Which is redder?}
\item \textbf{Which is bluer?}
\end{itemize}

If larger $\theta$ corresponds to redder colors on the chosen arc, then:
\begin{itemize}[leftmargin=2em]
\item in a \textbf{redder} context, the correct response is right when
$\theta_R > \theta_L$,
\item in a \textbf{bluer} context, the correct response is right when
$\theta_R < \theta_L$.
\end{itemize}

The task is therefore defined by a comparison of the two latent hue signals,
not by an absolute report of one stimulus alone.

\section{Two-Stage Per-Stimulus Noise Calibration}

The revised psychophysics paradigm uses the same 2AFC stimulus structure, but
changes the role of external noise. The goal is to use standard psychophysics
to estimate each observer's internal noise and then add subject-specific
external noise so that different observers have the same effective total
uncertainty in the task.

The paradigm has two stages:
\begin{enumerate}[leftmargin=2em]
\item \textbf{Experiment 1: baseline threshold measurement.} Measure a
baseline psychometric curve with no injected external theta noise and estimate
the subject's per-stimulus internal noise scale $\sigma_{\mathrm{int}}$.
\item \textbf{Experiment 2: matched-noise task.} Choose experimenter-defined
target per-stimulus total noise levels $\sigma_{\mathrm{target}}$ and compute
the external noise needed for each subject.
\end{enumerate}

\subsection{Per-Stimulus Versus Comparison Noise}

The sigmas in this calibration are defined per stimulus. For a left-right
trial, the observer is modeled as forming noisy latent-angle estimates
\[
\hat{\theta}_L
= \theta_L
+ \varepsilon_{L,\mathrm{int}}
+ \varepsilon_{L,\mathrm{ext}},
\qquad
\hat{\theta}_R
= \theta_R
+ \varepsilon_{R,\mathrm{int}}
+ \varepsilon_{R,\mathrm{ext}}.
\]

Under the default independent-noise model,
\[
\varepsilon_{L,\mathrm{int}}, \varepsilon_{R,\mathrm{int}}
\sim \mathcal{N}(0,\sigma_{\mathrm{int}}^2),
\qquad
\varepsilon_{L,\mathrm{ext}}, \varepsilon_{R,\mathrm{ext}}
\sim \mathcal{N}(0,\sigma_{\mathrm{ext}}^2),
\]
with all four terms independent.

The 2AFC decision is based on the perceived difference
\[
D = \hat{\theta}_R - \hat{\theta}_L.
\]

For a redder judgment,
\[
D
= \Delta\theta
+ (\varepsilon_{R,\mathrm{int}} - \varepsilon_{L,\mathrm{int}})
+ (\varepsilon_{R,\mathrm{ext}} - \varepsilon_{L,\mathrm{ext}}).
\]

Therefore,
\[
\operatorname{Var}(D)
= 2\sigma_{\mathrm{int}}^2 + 2\sigma_{\mathrm{ext}}^2
= 2\left(\sigma_{\mathrm{int}}^2 + \sigma_{\mathrm{ext}}^2\right).
\]

This factor of two appears because the subject compares two noisy stimulus
estimates. The calibration variables $\sigma_{\mathrm{int}}$,
$\sigma_{\mathrm{ext}}$, and $\sigma_{\mathrm{target}}$ remain
\textbf{per-stimulus} quantities.

\section{Experiment 1: Estimating Internal Noise}

The baseline experiment uses no injected external theta noise:
\[
\sigma_{\mathrm{ext}} = 0.
\]

For a redder judgment under the independent Gaussian 2AFC model,
\[
P(\text{choose right} \mid \Delta\theta)
= \Phi\!\left(
\frac{\Delta\theta}{\sqrt{2}\,\sigma_{\mathrm{int}}}
\right),
\]
where $\Phi$ is the standard normal cumulative distribution function.

The bluer judgment has the same form after flipping the sign of
$\Delta\theta$, so that positive effective evidence always favors a rightward
response.

If a psychometric threshold $\operatorname{JND}_p$ is defined as the value of
$|\Delta\theta|$ required to reach criterion performance $p$, then
\[
\operatorname{JND}_p
= \sqrt{2}\,\sigma_{\mathrm{int}}\,\Phi^{-1}(p).
\]

Solving for the per-stimulus internal noise gives
\[
\sigma_{\mathrm{int}}
= \frac{\operatorname{JND}_p}
{\sqrt{2}\,\Phi^{-1}(p)}.
\]

For example, if the threshold is defined at $p=0.75$, then
\[
\sigma_{\mathrm{int}}
= \frac{\operatorname{JND}_{0.75}}
{\sqrt{2}\,\Phi^{-1}(0.75)}.
\]

This is the key calibration step: the psychometric curve is fit to the
left-right comparison, but the reported internal noise is converted back to a
per-stimulus noise scale.

\section{Experiment 2: Matching Total Noise Across Subjects}

In the matched-noise experiment, the experimenter chooses a set of target
per-stimulus total noise levels,
\[
\sigma_{\mathrm{target},1},
\sigma_{\mathrm{target},2},
\dots,
\sigma_{\mathrm{target},K}.
\]

For each subject and each target level, external theta noise is chosen so that
\[
\sigma_{\mathrm{target}}^2
= \sigma_{\mathrm{int}}^2 + \sigma_{\mathrm{ext}}^2.
\]

Thus,
\[
\sigma_{\mathrm{ext}}
= \sqrt{
\sigma_{\mathrm{target}}^2 - \sigma_{\mathrm{int}}^2
}.
\]

This calculation is feasible only when
\[
\sigma_{\mathrm{target}} \ge \sigma_{\mathrm{int}}.
\]
If a subject's internal noise is already larger than a target level, then that
target cannot be reached by adding independent external noise. In practice,
target levels should be chosen at or above the largest internal-noise estimate
among the subjects intended to share a matched-noise condition.

Under this construction, two subjects with different baseline thresholds can
receive different external noise levels but have the same predicted total
per-stimulus uncertainty:
\[
\sigma_{\mathrm{total}}
= \sqrt{\sigma_{\mathrm{int}}^2 + \sigma_{\mathrm{ext}}^2}
= \sigma_{\mathrm{target}}.
\]

The predicted comparison-level standard deviation is then
\[
\sigma_D = \sqrt{2}\,\sigma_{\mathrm{target}},
\]
so the resulting psychometric curves should overlap when plotted against
$\Delta\theta$.

\subsection{SNR Interpretation}

The same calibration can be described in signal-to-noise-ratio language. For
a target hue separation $\Delta\theta$, the comparison-level SNR is
\[
\operatorname{SNR}
= \frac{|\Delta\theta|}{\sqrt{2}\,\sigma_{\mathrm{target}}}.
\]

The design rule is still total-noise matching. SNR is a useful way to report
the implied difficulty of a chosen $\Delta\theta$ under a chosen
$\sigma_{\mathrm{target}}$.

\section{Correlated Internal Noise Caveat}

The derivation above assumes that the left and right internal noises are
independent. If there are common internal fluctuations across the two
simultaneously viewed stimuli, then the internal noise terms may be correlated.

Let
\[
\operatorname{Corr}(
\varepsilon_{L,\mathrm{int}},
\varepsilon_{R,\mathrm{int}}
) = \rho_{\mathrm{int}}.
\]

Assuming both sides have the same per-stimulus internal variance
$\sigma_{\mathrm{int}}^2$, the internal contribution to the comparison
variance is
\[
\operatorname{Var}(
\varepsilon_{R,\mathrm{int}} - \varepsilon_{L,\mathrm{int}}
)
= 2\sigma_{\mathrm{int}}^2(1-\rho_{\mathrm{int}}).
\]

A standard 2AFC psychometric curve therefore identifies the effective
differential internal noise,
\[
\sigma_{\mathrm{int,diff}}
= \sigma_{\mathrm{int}}\sqrt{1-\rho_{\mathrm{int}}},
\]
not $\sigma_{\mathrm{int}}$ and $\rho_{\mathrm{int}}$ separately.

For the first version of the calibration analysis, the working assumption is
\[
\rho_{\mathrm{int}} = 0.
\]
If common internal fluctuations are plausible, the baseline calibration should
be interpreted more carefully as the per-stimulus noise component that
survives the left-right comparison.

Shared external-noise controls are useful here. If shared external noise
mostly cancels, then common-mode perturbations are not strongly affecting
discrimination. If shared-noise thresholds are not flat, then the calibration
should be reported as an effective differential-noise calibration rather than
a literal independent per-stimulus noise estimate. Estimating
$\rho_{\mathrm{int}}$ directly would require an additional design that can
separate common-mode and differential noise.

\section{Summary: Old Versus New Stimulus Definitions}

\subsection*{Old stimulus}

One stimulus is generated from
\[
X(u,v) = N(u,v) + \kappa e^{i\theta},
\]

followed by a pixel-wise angle readout and colormap mapping.

The key noise source is the correlated texture field $N(u,v)$.

\subsection*{New calibrated 2AFC design}

A trial contains a left-right pair with latent targets
\[
\theta_L^{\ast} = \theta_0 - \frac{\Delta\theta}{2},
\qquad
\theta_R^{\ast} = \theta_0 + \frac{\Delta\theta}{2},
\]

and then can add trial-level theta noise before rendering each side through
the same generator.

The key new manipulation is not a change in the pixel generator itself, but a
change in the \textbf{latent trial structure}: baseline 2AFC behavior is used
to estimate per-stimulus internal noise, and external theta noise is then
chosen subject by subject to reach experimenter-defined total noise levels.

\end{document}
